\section{Dekomposisi Hamiltonian: Dari Matriks Global ke String Pauli}\label{appendix-e}

\hypertarget{urgensi-eksplorasi-masalah-raksasa-di-ruang-hilbert}{%
\subsection{Urgensi Eksplorasi: Masalah ``Raksasa'' di Ruang Hilbert}\label{urgensi-eksplorasi-masalah-raksasa-di-ruang-hilbert}}

Sistem \(n\)-\emph{qubit} memiliki ruang status sebesar \(2^n \times 2^n\). Untuk 10 aset saja, matriks Hamiltonian berukuran \(1024 \times 1024\). Komputer kuantum tidak dapat ``membaca'' atau memproses matriks sebesar ini secara utuh. Hardware kuantum hanya dapat melakukan pengukuran pada basis tertentu (biasanya basis Z). Kita butuh cara untuk memecah informasi global ini menjadi potongan-potongan lokal yang bisa diukur secara paralel.

\begin{center}\rule{0.5\linewidth}{0.5pt}\end{center}

\hypertarget{aksioma-intuisi-matriks-pauli-sebagai-atom-informasi}{%
\subsection{Aksioma \& Intuisi: Matriks Pauli sebagai ``Atom'' Informasi}\label{aksioma-intuisi-matriks-pauli-sebagai-atom-informasi}}

Setiap matriks Hermitian (energi) dapat dibangun dari 4 ``batu bata'' dasar yang disebut \textbf{Matriks Pauli}:

\[ I = \begin{pmatrix} 1 & 0 \\ 0 & 1 \end{pmatrix}, X = \begin{pmatrix} 0 & 1 \\ 1 & 0 \end{pmatrix}, Y = \begin{pmatrix} 0 & -i \\ i & 0 \end{pmatrix}, Z = \begin{pmatrix} 1 & 0 \\ 0 & -1 \end{pmatrix} \]

\textbf{Aksioma:} Matriks Pauli \(\{I, X, Y, Z\}\) membentuk basis ortogonal yang lengkap bagi ruang matriks \(2 \times 2\). Artinya, sembarang matriks energi 1-\emph{qubit} pasti bisa dinyatakan sebagai kombinasi linear dari keempat matriks ini dengan koefisien riil.

\begin{center}\rule{0.5\linewidth}{0.5pt}\end{center}

\hypertarget{reduksionisme-kasus-minimal-1-qubit}{%
\subsection{Reduksionisme: Kasus Minimal (1-Qubit)}\label{reduksionisme-kasus-minimal-1-qubit}}

Mari kita bedah sebuah Hamiltonian 1-\emph{qubit} umum \(\hat{H}\): \[ \hat{H} = c_0 I + c_x X + c_y Y + c_z Z \qquad (1) \]

\hypertarget{a.-mencari-koefisien-proyeksi-trace}{%
\subsubsection{Mencari Koefisien (Proyeksi Trace)}\label{a.-mencari-koefisien-proyeksi-trace}}

Kita memproyeksikan \(\hat{H}\) ke masing-masing basis Pauli menggunakan operasi \textbf{Trace} (\(\text{Tr}\)): \[ c_k = \frac{1}{2} \text{Tr}(\sigma_k \hat{H}) \qquad (2) \]

\begin{quote}
\textbf{Visualisasi (2): Perhitungan Matriks (Mencari \(c_z\))} Misalkan kita punya Hamiltonian dengan energi state \(|0\rangle\) sebesar 3 dan state \(|1\rangle\) sebesar 1, maka \(\hat{H} = \begin{pmatrix} 3 & 0 \\ 0 & 1 \end{pmatrix}\). Mari kita cari komponen \(Z\): \[ 
\begin{split}
c_z &= \frac{1}{2} \text{Tr} \left[ \begin{pmatrix} 1 & 0 \\ 0 & -1 \end{pmatrix} \begin{pmatrix} 3 & 0 \\ 0 & 1 \end{pmatrix} \right] \\\\
c_z &= \frac{1}{2} \text{Tr} \begin{pmatrix} 3 & 0 \\ 0 & -1 \end{pmatrix} = \frac{1}{2} (3 + (-1)) = 1
\end{split}
\] Artinya, energi ini mengandung ``gaya'' sebesar \(1\) unit ke arah sumbu Z.
\end{quote}

\begin{center}\rule{0.5\linewidth}{0.5pt}\end{center}

\hypertarget{jembatan-formalisme-logika-banyak-qubit-many-body}{%
\subsection{Jembatan Formalisme \& Logika: Banyak Qubit (Many-Body)}\label{jembatan-formalisme-logika-banyak-qubit-many-body}}

Untuk sistem \(n\)-\emph{qubit}, kita menggunakan \textbf{Produk Tensor} (\(\otimes\)) untuk membangun basis yang lebih besar.

\textbf{Jembatan Logika:} Jika 1 \emph{qubit} punya 4 basis, maka \(n\) \emph{qubit} punya \(4^n\) kombinasi basis (Pauli Strings). Hamiltonian total didekomposisi menjadi jumlahan string: \[ \hat{H} = \sum_{k=1}^{4^n} c_k P_k \qquad (3) \] Di mana \(P_k\) adalah string seperti \(Z \otimes I \otimes X \otimes \dots\).

\begin{center}\rule{0.5\linewidth}{0.5pt}\end{center}

\hypertarget{derivasi-scratchpad-indeks-persamaan-linearitas-ekspektasi}{%
\subsection{Derivasi ``Scratchpad'' \& Indeks Persamaan: Linearitas Ekspektasi}\label{derivasi-scratchpad-indeks-persamaan-linearitas-ekspektasi}}

Mengapa dekomposisi ini menjadi kunci VQE? Karena sifat linearitas operator dalam mekanika kuantum: \[ \langle \hat{H} \rangle = \langle \psi | \left( \sum_k c_k P_k \right) | \psi \rangle = \sum_k c_k \langle \psi | P_k | \psi \rangle \qquad (4) \]

\begin{quote}
\textbf{Visualisasi (4): Aliran Data} Bayangkan kita ingin mengukur energi total: 1. \textbf{Klasik:} Komputer menyimpan daftar koefisien \(\{c_k\}\) (misal: \(c_1=0.5, c_2=-1.2\)). 2. \textbf{Kuantum:} Hardware mengukur rata-rata string \(\langle P_k \rangle\) (misal: \(\langle Z_1 Z_2 \rangle = 0.8\)). 3. \textbf{Hasil:} Kita kalikan secara klasik \((0.5 \times 0.8 + \dots)\) untuk mendapatkan total energi.

Ini memungkinkan kita menghitung energi raksasa hanya dengan melakukan banyak pengukuran kecil yang sederhana.
\end{quote}

\begin{center}\rule{0.5\linewidth}{0.5pt}\end{center}

\hypertarget{verifikasi-parameter-portofolio-markowitz-ising}{%
\subsection{Verifikasi \& Parameter: Portofolio Markowitz-Ising}\label{verifikasi-parameter-portofolio-markowitz-ising}}

Dalam masalah optimasi portofolio, dekomposisi ini sangat efisien karena Hamiltonian Ising hanya memiliki dua jenis ``Lego'':

\begin{longtable}[]{@{}
  >{\raggedright\arraybackslash}p{(\columnwidth - 6\tabcolsep) * \real{0.2500}}
  >{\raggedright\arraybackslash}p{(\columnwidth - 6\tabcolsep) * \real{0.2500}}
  >{\raggedright\arraybackslash}p{(\columnwidth - 6\tabcolsep) * \real{0.2500}}
  >{\raggedright\arraybackslash}p{(\columnwidth - 6\tabcolsep) * \real{0.2500}}@{}}
\toprule\noalign{}
\begin{minipage}[b]{\linewidth}\raggedright
Jenis Suku
\end{minipage} & \begin{minipage}[b]{\linewidth}\raggedright
Komponen Fisik
\end{minipage} & \begin{minipage}[b]{\linewidth}\raggedright
Pauli String (\(P_k\))
\end{minipage} & \begin{minipage}[b]{\linewidth}\raggedright
Makna Ekonomi
\end{minipage} \\
\midrule\noalign{}
\endhead
\bottomrule\noalign{}
\endlastfoot
\textbf{Suku Linear} & \(h_i \hat{Z}_i\) & \(I \otimes \dots \otimes Z \otimes \dots \otimes I\) & Risiko/Return individu aset \(i\). \\
\textbf{Suku Interaksi} & \(J_{ij} \hat{Z}_i \hat{Z}_j\) & \(I \otimes \dots \otimes Z \otimes \dots \otimes Z \otimes \dots \otimes I\) & Kovarians (hubungan gerak) aset \(i\) dan \(j\). \\
\end{longtable}

\begin{center}\rule{0.5\linewidth}{0.5pt}\end{center}

\hypertarget{analogi-physical-insight-spektroskopi-operator}{%
\subsection{Analogi ``Physical Insight'': Spektroskopi Operator}\label{analogi-physical-insight-spektroskopi-operator}}

Bayangkan Hamiltonian sebagai \textbf{cahaya putih} yang misterius. - Dekomposisi Hamiltonian bertindak seperti \textbf{Prisma} yang memecah cahaya tersebut menjadi warna-warna dasar (Pauli Strings). - Komputer kuantum adalah \textbf{Detektor} yang hanya bisa melihat satu warna dalam satu waktu. - Dengan mengukur intensitas setiap warna dasar, kita bisa merekonstruksi spektrum energi total dari cahaya tersebut tanpa harus menelan seluruh matari pembentuknya.

\begin{center}\rule{0.5\linewidth}{0.5pt}\end{center}

\hypertarget{kesimpulan}{%
\subsection{Kesimpulan}\label{kesimpulan-e}}

Dekomposisi Hamiltonian mengubah masalah \textbf{eksponensial} (matriks raksasa) menjadi jumlahan \textbf{polinomial} dari pengukuran-pengukuran lokal yang praktis. Tanpa teknik ini, prinsip variasi (Rayleigh-Ritz) tidak akan pernah bisa diimplementasikan pada \emph{hardware} kuantum era \emph{NISQ}.
